
        \documentclass[fleqn]{article}      
        \usepackage{amsmath}
        \usepackage{fontspec}
        \usepackage{kotex} % 한국어 지원  

        \begin{document}      
        ```latex
 쉬운 문제들
\textbf{쉬운 문제들}
\noindent 1. 다음 행렬 A의 차원은 무엇인가?
\begin{align*}
A = \begin{pmatrix}
1 & 2 \\
3 & 4
\end{pmatrix}
\end{align*}
\\\\

\noindent 2. 두 행렬 B와 C가 다음과 같을 때, B + C의 결과는 무엇인가?
\begin{align*}
B = \begin{pmatrix}
1 & 0 \\
0 & 1
\end{pmatrix}, \quad C = \begin{pmatrix}
2 & 2 \\
2 & 2
\end{pmatrix}
\end{align*}
\\\\

\noindent 3. 단위 행렬 I의 정의는 무엇인가? 
\\\\

\noindent 4. 행렬 D의 주대각선의 합을 구하시오.
\begin{align*}
D = \begin{pmatrix}
5 & 1 \\
2 & 3
\end{pmatrix}
\end{align*}
\\\\

\noindent 5. 행렬 E의 전치 행렬 E^T는 무엇인가?
\begin{align*}
E = \begin{pmatrix}
7 & 8 \\
9 & 10
\end{pmatrix}
\end{align*}
\\\\

 중간 난이도 문제들
\textbf{중간 난이도 문제들}
\noindent 6. 행렬 F의 곱을 구하시오.
\begin{align*}
F = \begin{pmatrix}
1 & 2 \\
3 & 4
\end{pmatrix}, \quad G = \begin{pmatrix}
5 & 6 \\
7 & 8
\end{pmatrix}
\end{align*}
\\\\

\noindent 7. 다음 행렬의 역행렬을 구하시오.
\begin{align*}
H = \begin{pmatrix}
1 & 2 \\
3 & 4
\end{pmatrix}
\end{align*}
\\\\

\noindent 8. 다음 행렬의 랭크를 구하시오.
\begin{align*}
I = \begin{pmatrix}
1 & 2 & 3 \\
4 & 5 & 6 \\
7 & 8 & 9
\end{pmatrix}
\end{align*}
\\\\

\noindent 9. 행렬 J의 고유값을 구하시오.
\begin{align*}
J = \begin{pmatrix}
2 & 1 \\
1 & 2
\end{pmatrix}
\end{align*}
\\\\

\noindent 10. 두 행렬 K와 L의 내적을 계산하시오.
\begin{align*}
K = \begin{pmatrix}
1 & 2 \\
3 & 4
\end{pmatrix}, \quad L = \begin{pmatrix}
5 & 6 \\
7 & 8
\end{pmatrix}
\end{align*}
\\\\

 어려운 문제들
\textbf{어려운 문제들}
\noindent 11. 주어진 행렬 M의 행렬식(det)을 구하시오.
\begin{align*}
M = \begin{pmatrix}
3 & 2 & 1 \\
1 & 0 & 2 \\
0 & 1 & 3
\end{pmatrix}
\end{align*}
\\\\

\noindent 12. 다음 세트의 행렬 N, O, P의 곱을 계산하시오.
\begin{align*}
N = \begin{pmatrix}
1 & 0 \\
0 & 1
\end{pmatrix}, \quad O = \begin{pmatrix}
2 & 1 \\
1 & 2
\end{pmatrix}, \quad P = \begin{pmatrix}
3 & 0 \\
0 & 3
\end{pmatrix}
\end{align*}
\\\\

\noindent 13. 행렬 Q의 고유벡터를 구하시오.
\begin{align*}
Q = \begin{pmatrix}
4 & 1 \\
2 & 3
\end{pmatrix}
\end{align*}
\\\\

\noindent 14. 다음 선형방정식 Ax = b를 풀어라.
\begin{align*}
A = \begin{pmatrix}
1 & 2 \\
3 & 4
\end{pmatrix}, \quad b = \begin{pmatrix}
5 \\
6
\end{pmatrix}
\end{align*}
\\\\

\noindent 15. 다음 행렬 R의 고유값과 고유벡터를 구하시오.
\begin{align*}
R = \begin{pmatrix}
0 & 1 \\
-1 & 0
\end{pmatrix}
\end{align*}
\\\\

 답안
\textbf{답안}
1. 2x2 행렬 \\\\\\
2. \begin{pmatrix} 3 & 2 \\ 2 & 3 \end{pmatrix} \\\\\\
3. 모든 주대각선의 원소가 1이고 나머지가 0인 정사각형 행렬 \\\\\\
4. 8 \\\\\\
5. \begin{pmatrix} 7 & 9 \\ 8 & 10 \end{pmatrix} \\\\\\
6. \begin{pmatrix} 19 & 22 \\ 43 & 50 \end{pmatrix} \\\\\\
7. \begin{pmatrix} -2 & 1 \\ 1.5 & -0.5 \end{pmatrix} \\\\\\
8. 2 \\\\\\
9. 3 ± √2 \\\\\\
10. 70 \\\\\\
11. 1 \\\\\\
12. \begin{pmatrix} 6 & 3 \\ 3 & 6 \end{pmatrix} \\\\\\
13. \begin{pmatrix} 1 \\ -0.5 \end{pmatrix}, \begin{pmatrix} -1 \\ 2 \end{pmatrix} \\\\\\
14. \begin{pmatrix} 2 \\ 1 \end{pmatrix} \\\\\\
15. 고유값: {1, -1}, 고유벡터: \begin{pmatrix} 1 \\ 1 \end{pmatrix}, \begin{pmatrix} -1 \\ 1 \end{pmatrix} \\\\\\
```      
        \end{document} 
    